\subsection{Circle and ellipse: from one scale to two}

A circle centred at the origin is described by
\[
\boxed{x^2+y^2=R^2}.
\]
Its equation treats every direction equally. The parameterisation
\[
\boxed{x=R\cos t,\qquad y=R\sin t}
\]
makes this symmetry visible: one angle determines one point on the circle.

An ellipse introduces two perpendicular scales:
\[
\boxed{\frac{x^2}{a^2}+\frac{y^2}{b^2}=1},
\qquad a\geq b>0.
\]
It is parameterised by
\[
\boxed{x=a\cos t,\qquad y=b\sin t}.
\]
The same circular parameter is stretched by different factors in the two
coordinate directions.

\begin{center}
\begin{tikzpicture}[scale=0.82,>=stealth]
    \begin{scope}[xshift=-4.2cm]
        \draw[axis,->] (-3.0,0)--(3.2,0) node[right] {$x$};
        \draw[axis,->] (0,-2.8)--(0,3.0) node[above] {$y$};
        \draw[subset] (0,0) circle (2.0);
        \draw[blue!70!black,line width=1pt,->] (0,0)--(1.53,1.29)
            node[midway,above left] {$R$};
        \node at (0,-3.4) {$x^2+y^2=R^2$};
    \end{scope}
    \begin{scope}[xshift=4.2cm]
        \draw[axis,->] (-4.0,0)--(4.2,0) node[right] {$x$};
        \draw[axis,->] (0,-2.8)--(0,3.0) node[above] {$y$};
        \draw[subset] (0,0) ellipse (3.1 and 1.8);
        \draw[blue!70!black,line width=1pt] (0,0)--(3.1,0)
            node[midway,below] {$a$};
        \draw[orange!85!black,line width=1pt] (0,0)--(0,1.8)
            node[midway,left] {$b$};
        \node at (0,-3.4) {$x^2/a^2+y^2/b^2=1$};
    \end{scope}
\end{tikzpicture}
\end{center}

The ellipse also has a second geometric description. Define
\[
c=\sqrt{a^2-b^2}.
\]
Its foci are $F_1=(-c,0)$ and $F_2=(c,0)$, and every point $P$ on the ellipse
satisfies
\[
\boxed{d(P,F_1)+d(P,F_2)=2a}.
\]

\begin{center}
\begin{tikzpicture}[scale=0.9,>=stealth]
    \draw[axis,->] (-4.5,0)--(4.7,0);
    \draw[axis,->] (0,-2.7)--(0,2.9);
    \draw[subset] (0,0) ellipse (3.5 and 2.0);
    \coordinate (F1) at (-2.87,0);
    \coordinate (F2) at (2.87,0);
    \coordinate (P) at (1.4,1.83);
    \draw[blue!70!black,line width=1.1pt] (F1)--(P)
        node[midway,above left] {$r_1$};
    \draw[orange!85!black,line width=1.1pt] (P)--(F2)
        node[midway,above right] {$r_2$};
    \fill[geometricSubsetColor] (F1) circle (2.2pt) node[below] {$F_1$};
    \fill[geometricSubsetColor] (F2) circle (2.2pt) node[below] {$F_2$};
    \fill[geometricSubsetColor] (P) circle (2.2pt) node[above] {$P$};
    \node at (0,-3.25) {$r_1+r_2=2a$};
\end{tikzpicture}
\end{center}

\begin{interpretationbox}[Two useful descriptions]
The centre-based equation exposes symmetry and semiaxes. The focal description
exposes distances from the points that later become natural locations of a
central force source.
\end{interpretationbox}
